\documentclass[tikz,border=6pt]{standalone}
\usepackage{peabra}
\begin{document}
\begin{tikzpicture}[x=1cm,y=1cm, every node/.append style={font=\scriptsize}]
  \foreach \y/\c/\szint/\leir/\hiba/\pl in {
     0/peKek/{készségalapú}/{begyakorolt, automatikus mozdulatok; kevés tudatos figyelem}/{\textbf{csúszás:} rossz kivitelezés\\ \textbf{kihagyás:} egy lépés elmarad}/{rossz kapcsolót nyom meg; elfelejti visszaállítani a kiiktatást},
    -1.6/peZold/{szabályalapú}/{ismerős helyzet, betanult szabály: ha X, akkor Y}/{\textbf{szabályalapú tévedés:} jó szabály rossz helyzetben, vagy rossz szabály}/{a szokásos riasztásnak véli a valódi veszélyt},
    -3.2/peKiemel/{tudásalapú}/{új, ismeretlen helyzet; tudatos következtetés, mentális modell}/{\textbf{tudásalapú tévedés:} hiányos mentális modell, téves diagnózis}/{nem ismeri fel az átesést, mert a figyelmeztetés ellentmondásos}} {
    \node[draw=\c, fill=\c!15, rounded corners=2pt, text width=26mm, minimum height=13mm, align=center, font=\small\bfseries] (s) at (0,\y) {\szint};
    \node[text width=40mm, align=left, anchor=west] at (1.7,\y) {\leir};
    \node[draw=\c, fill=white, rounded corners=2pt, text width=44mm, minimum height=13mm, align=left, anchor=west, inner sep=3pt] at (6.2,\y) {\hiba};
    \node[text width=38mm, align=left, anchor=west, text=peSzurke, font=\tiny] at (11.0,\y) {\pl};
  }
  \node[font=\small\bfseries, text=peSotet] at (0,1.1) {Rasmussen: viselkedési szint};
  \node[font=\small\bfseries, text=peSotet, anchor=west] at (6.2,1.1) {Reason: hibatípus};
  \node[font=\small\bfseries, text=peSotet, anchor=west] at (11.0,1.1) {példa};
  \draw[nyil, draw=peSzurke] (-1.6,0.6) -- node[left, align=right, font=\tiny, text width=14mm]{nő a tudatos figyelem és a hiba valószínűsége} (-1.6,-3.8);
  \node[pirosdoboz=148mm, minimum height=8mm, align=left, anchor=north west, inner sep=3pt] at (-1.3,-4.3) {\textbf{szabályszegés:} szándékos eltérés az előírástól (rutinszerű, helyzeti vagy kivételes); nem tévedés, ezért más eszközök kezelik: az előírások gyakorlatiassága, a szervezeti kultúra, a munkakörülmények};
\end{tikzpicture}
\end{document}
